\documentclass[11pt]{article}

\usepackage[margin=1in]{geometry}
\usepackage{amsmath,amssymb,amsthm,mathtools}
\usepackage{booktabs}
\usepackage[hidelinks]{hyperref}

\newtheorem{theorem}{Theorem}[section]
\newtheorem{lemma}[theorem]{Lemma}
\newtheorem{proposition}[theorem]{Proposition}
\newtheorem{corollary}[theorem]{Corollary}
\theoremstyle{remark}
\newtheorem{remark}[theorem]{Remark}

\newcommand{\splus}{s^+}
\newcommand{\sminus}{s^-}
\newcommand{\sig}{\sigma}
\newcommand{\D}{D}
\newcommand{\tr}{\operatorname{tr}}
\newcommand{\T}{\mathsf T}
\newcommand{\Pp}{\mathsf P}

\title{Positive Square Energy for Cacti with Triangles and Two Pentagons}
\author{AI-generated manuscript; no human author is claimed}
\date{30 July 2026}

\begin{document}
\maketitle

\begin{abstract}
Let $G$ be a connected cactus whose complete cyclic blocks are $r\geq1$
triangles and exactly two pentagons.  Shared cut vertices, bridge connectors of
arbitrary length and branching, and arbitrary finite trees attached at
arbitrary vertices are all allowed.  We prove
\[
  \splus(G)>|V(G)|.
\]
The proof applies maximum-triangle Voronoi to the original graph before any
bridge cut.  Every resulting territory has triangle-packing number one.  We
give a self-contained Sachs--Coulson proof of the required one-pentagon local
estimate with unrestricted incidence and ports.  A well-founded actual-bridge
induction reduces a territory retaining both pentagons to one shared-cut
cluster.  Its incidence tree has exactly seven structural forms; common-cut
and rooted two-pentagon hinge theorems handle three forms, and explicit
physical interval ownership handles the others.  Every cut, fragment,
connector remnant, and attached tree receives exactly one owner.
\end{abstract}

\section{Statement, notation, and established packet inputs}

All graphs are finite, simple, and undirected.  For the adjacency eigenvalues
$\lambda_1,\ldots,\lambda_n$ of a graph $H$, put
\[
 \splus(H)=\sum_{\lambda_j>0}\lambda_j^2,
 \qquad
 \sminus(H)=\sum_{\lambda_j<0}\lambda_j^2,
 \qquad
 \sig(H)=\splus(H)-|V(H)|,
\]
and write $\D(H)=\splus(H)-\sminus(H)$.  A cactus is a graph whose blocks are
bridges, isolated vertices, or cycles.  Throughout, $\T=C_3$, $\Pp=C_5$, and
\[
 \delta=\sqrt5-2.
\]
A ``complete cycle'' in an induced subgraph means a cyclic block retained in
its entirety; a proper fragment of a cycle is only a forest.

\begin{theorem}\label{thm:main}
Let $G$ be a connected cactus whose complete cyclic blocks are $r\geq1$
triangles and exactly two copies of $C_5$.  The cyclic blocks may meet at cut
vertices in arbitrary cactus incidence, distinct cyclic clusters may be joined
by arbitrary bridge trees, and arbitrary finite trees may be attached at every
vertex.  If $n=|V(G)|$, then
\[
 \boxed{\splus(G)>n}.
\]
\end{theorem}

We use the following proved packet results.  They are stated in precisely the
tree-uniform forms needed here; their sources are listed at the end.

\begin{lemma}[Packet inputs]\label{lem:inputs}
The following assertions hold with arbitrary finite tree attachments.
\begin{enumerate}
\item\label{in:super} If $V(H)=V_1\dot\cup\cdots\dot\cup V_k$, then
\[
 \splus(H)\geq\sum_j\splus(H[V_j]),
 \qquad
 \sig(H)\geq\sum_j\sig(H[V_j]).
\]
\item\label{in:tri} If a connected cactus has $b\geq1$ triangular cyclic
blocks, no other complete cycles, and no two vertex-disjoint triangles, then
\[
 \sig(H)>b-1\geq0.
\]
\item\label{in:pents} A connected unicyclic pentagonal cactus has
$\sig(H)\geq-\delta$.  A connected cactus whose only complete cycles are two
pentagons has $\sig(H)\geq0$.
\item\label{in:common} If $b\geq1$ triangles and two pentagons all contain one
common cut vertex, then
\[
 \sig(H)>b+1-\frac4{3\sqrt{13}}>0.
\]
\item\label{in:small} Every connected tricyclic cactus has nonnegative
surplus, and every $\T\Pp\Pp$ cactus has positive surplus.  The rank-four and
rank-five packets $\T^2\Pp\Pp$ and $\T^3\Pp\Pp$ also have positive surplus.
\item\label{in:hinge} Suppose two pentagons share a cut $x$, and a nonempty
connected triangular cactus enters at one vertex $y$ of the first pentagon
through one interface, where $y=x$ is allowed.  Then the resulting rooted
two-pentagon hinge has positive surplus.
\end{enumerate}
\end{lemma}

For completeness, item~\ref{in:super} follows from the variational formula
\[
 \tr(A_+^2)=\max_{X\succeq0}
 \bigl(2\tr(AX)-\tr(X^2)\bigr).
\]
Order the vertices by the partition and test with the positive spectral part
of the block-diagonal matrix having diagonal blocks $A(H[V_j])$.  Off-block
entries make no trace contribution, proving the first inequality; subtracting
$\sum_j|V_j|=|V(H)|$ proves the second.  Items
\ref{in:tri}--\ref{in:hinge} are the common-cut, hinge, and low-rank companion
theorems cited in Section~\ref{sec:dependencies}.  The new estimate required
for this proof is Lemma~\ref{lem:L}; unlike an earlier one-interface packet,
it permits any number of pentagonal ports.

\section{Voronoi first}\label{sec:voronoi}

Choose a maximum-cardinality family
$\T_1,\ldots,\T_m$ of pairwise vertex-disjoint triangles in the original
graph $G$.  Assign each vertex $v$ to the index minimizing lexicographically
\[
 \bigl(d_G(v,V(\T_i)),i\bigr),                              \tag{2.1}
\]
and let $G_i$ be induced by the vertices assigned to $i$.

\begin{lemma}[Maximum-triangle Voronoi territories]\label{lem:voronoi}
The sets $V(G_i)$ form an exhaustive induced vertex partition.  Every $G_i$
is connected, contains the whole selected triangle $\T_i$, and contains no
two vertex-disjoint retained triangles.
\end{lemma}

\begin{proof}
Certainly every vertex has exactly one owner.  Every vertex of $\T_i$ has
distance zero from $\T_i$; disjointness of the selected family makes $i$ its
unique zero-distance owner, so $\T_i$ is retained wholly.

Here is the owner detail needed for connectivity.  Suppose $v$ has owner $i$
and $u$ is the predecessor of $v$ on a shortest path from $v$ to $\T_i$.
Then $d_i(u)=d_i(v)-1$.  For every $j$, one-Lipschitzness gives
$d_j(u)\geq d_j(v)-1$.  Since $(d_i(v),i)\leq(d_j(v),j)$ lexicographically,
it follows that
\[
 (d_i(u),i)=(d_i(v)-1,i)\leq(d_j(u),j).
\]
Thus $u$ has the same owner as $v$.  Iterating reaches $\T_i$ inside $G_i$,
so $G_i$ is connected.  This also clarifies that fixed priority is used only
to resolve equal distances consistently; it does not reassign an anchor after
the fact.

If $G_i$ contained vertex-disjoint triangles $D,E$, then
\[
 D,E,\{\T_j:j\ne i\}
\]
would be a triangle packing larger than the chosen one.  The two triangles
$D,E$ replace $\T_i$; they need not be disjoint from it.  This is why the
packing must be maximum, not merely maximal.
\end{proof}

An induced subgraph of a cactus creates no new cycle.  Therefore a triangle or
pentagon split among owners contributes only forest fragments; in particular,
a proper pentagon fragment is a path, never a new $\Pp$ packet.  A fragment may
lie on the minimal hull between retained cycles or hang off it, and both roles
are allowed as forest structure in the packet statements.  Every original
off-hull tree component has one actual anchor and follows the owner of that
anchor.  Thus no physical vertex is discarded, no fragment is promoted to a
complete cycle, and no attached tree is duplicated.

Classify $G_i$ by the number of complete pentagons it retains.
\begin{itemize}
\item A $0\Pp$ territory retains $b\geq1$ triangles and has triangle-packing
number one, so Lemma~\ref{lem:inputs}\ref{in:tri} makes it strict.
\item A $1\Pp$ territory retains $a\geq1$ packing-one triangles and is
positive by Lemma~\ref{lem:L} below, irrespective of ports and bridge
incidence.
\item A $2\Pp$ territory retains both pentagons and $a\geq1$ packing-one
triangles.  Sections~\ref{sec:bridges}--\ref{sec:H} prove it positive.
\end{itemize}
The territories are induced, disjoint, and exhaustive.  Consequently these
three local assertions, followed by Lemma~\ref{lem:inputs}\ref{in:super}, prove
the main theorem.

\section{Lemma L: unrestricted packing-one one-pentagon packets}

For a graph $F$ and positive vertex activities
$\alpha=(\alpha_v)_{v\in V(F)}$, define the signless matching partition
\[
 Z_F(\alpha)=\sum_{M\text{ matching of }F}
       \prod_{v\text{ unmatched by }M}\alpha_v,
 \qquad
 Z_F(t)=\sum_M t^{|V(F)|-2|M|}.
\]
Empty graphs have matching partition one.  Also normalize the characteristic
polynomial by
\[
 \Psi_F(t)=i^{-|V(F)|}\det(itI-A(F))
           =\prod_{j=1}^{|V(F)|}(t+i\lambda_j),\qquad t>0.
\]
Its continuous phase, tending to zero at infinity, is
\[
 \Theta_F(t)=\sum_j\arctan\frac{\lambda_j}{t}.
\]

\begin{lemma}[Signed Coulson identity]\label{lem:coulson}
For every finite graph $F$,
\[
 \D(F)=-\frac4\pi\int_0^\infty t\Theta_F(t)\,dt.          \tag{3.1}
\]
\end{lemma}

\begin{proof}
For real $\lambda$,
\[
 \lambda|\lambda|=\frac2\pi\int_0^\infty
       \frac{\lambda^3}{\lambda^2+t^2}\,dt.
\]
Moreover $\Theta_F'(t)=-\sum_j\lambda_j/(\lambda_j^2+t^2)$, and
$\sum_j\lambda_j=\tr A(F)=0$ gives
\[
 \sum_j\frac{\lambda_j^3}{\lambda_j^2+t^2}
 =t^2\Theta_F'(t).
\]
Sum the scalar identity and integrate by parts.  The boundary term vanishes:
$t^2\Theta_F(t)\to0$ at zero, and the arctangent expansion gives
$\Theta_F(t)=O(t^{-3})$ at infinity.  This also gives convergence and fixes
the sign in~(3.1).
\end{proof}

\begin{lemma}[Lemma L]\label{lem:L}
Let $H$ be a connected cactus whose complete cyclic blocks are one pentagon
$Q=C_5$ and $a\geq1$ triangles.  Suppose no two retained triangles are
vertex-disjoint.  Shared cuts, bridges, the number of occupied vertices of
$Q$, and arbitrary finite tree attachments are unrestricted.  Then
\[
 \D(H)>-2\delta,
 \qquad
 \sig(H)>a-\delta>0.                                     \tag{L}
\]
\end{lemma}

\begin{proof}
Let the \emph{spine} $S_H$ be the union of all complete cyclic blocks and the
unique minimal paths joining them.  Every component outside $S_H$ is a tree
with one spine attachment; two attachments would create another cycle.

Orient an off-spine branch toward the spine.  For an oriented edge $u\to v$,
let $T_{u\to v}$ be the branch below $u$, including $u$.  Splitting matchings
according to whether $u$ is unmatched or is matched to one child gives the
exact message recursion
\[
 M_{u\to v}(t)=\frac{Z_{T_{u\to v}}(t)}
                     {Z_{T_{u\to v}-u}(t)}
 =t+\sum_{w\text{ child of }u}\frac1{M_{w\to u}(t)}\geq t. \tag{3.2}
\]
Starting with leaves and extracting the denominators in~(3.2) produces the
same positive real factor $K(t)$ in every Sachs carrier and replaces the
unmatched-vertex activity at each spine vertex by
\[
 \alpha_v(t)=t+y_v(t),\qquad y_v(t)\geq0.                 \tag{3.3}
\]
The common-factor statement remains true when a Sachs cycle deletes a spine
vertex: the branch below it then contributes its already extracted full tree
factor, while an available root contributes the corresponding reciprocal
message.  Thus~(3.2)--(3.3) are an exact matching decomposition, not an
attachment monotonicity assertion.

We now derive the grouped Sachs formula.  For a collection $\mathcal C$ of
pairwise vertex-disjoint cycles, normalized Sachs expansion gives
\[
 \frac{\Psi_H(t)}{K(t)}
 =\sum_{\mathcal C}\prod_{C\in\mathcal C}(-2i^{-|C|})
     Z_{S_H-V(\mathcal C)}(\alpha).                       \tag{3.4}
\]
Indeed, after the cycle components are fixed, all remaining elementary
components form a matching; substitution at $it$ converts its alternating
matching polynomial to the signless partition.  A triangle has multiplier
$-2i$, while $Q=C_5$ has multiplier $+2i$.  Since no two triangles are
disjoint, a Sachs collection contains at most one triangle.  A disjoint
$Q$--triangle pair has multiplier $(2i)(-2i)=4$.  Therefore, exactly,
\[
 \frac{\Psi_H(t)}{K(t)}=R+2i(B-A),                        \tag{3.5}
\]
where, writing $S=V(Q)$,
\begin{align*}
 B&=Z_{S_H-S}(\alpha)>0,\\
 A&=\sum_T Z_{S_H-V(T)}(\alpha)>0,\\
 R&=Z_{S_H}(\alpha)
   +4\sum_{T\cap Q=\varnothing}
       Z_{S_H-(S\cup V(T))}(\alpha)>0.                   \tag{3.6}
\end{align*}
The sums run over retained triangles, and the last sum is zero if there is no
triangle disjoint from $Q$.

Nothing in~(3.5)--(3.6) assumes one port on $Q$.  Retain every shared cut of
$Q$ in $S$.  Every edge from such a cut into another block, and every first
edge of a bridge connector leaving $Q$, is then a boundary edge between $S$
and $S_H-S$.  Partition matchings counted by $Z_{S_H}(\alpha)$ according to
whether they use any such boundary edge.  Matchings using none factor, while
all other weights are nonnegative.  Hence
\[
 Z_{S_H}(\alpha)=Z_Q(\alpha|_S)B+E,\qquad E\geq0.          \tag{3.7}
\]
Because each activity is at least $t$ and matching partitions have
nonnegative coefficients,
\[
 Z_Q(\alpha|_S)=Z_5(t)+L,\qquad L\geq0,                   \tag{3.8}
\]
where $Z_5(t)$ is the bare $C_5$ matching partition.  Combining
(3.5)--(3.8) in the displayed notation gives
\begin{align*}
 R-Z_5(t)(B-A)
 &=E+LB+Z_5(t)A\\
 &\quad+4\sum_{T\cap Q=\varnothing}
       Z_{S_H-(S\cup V(T))}(\alpha)>0.                   \tag{3.9}
\end{align*}
The inequality is strict because $A>0$ and $Z_5(t)>0$.

Since $K,R>0$, the continuous phase is the principal value
\[
 \Theta_H(t)=\arctan\frac{2(B-A)}R.
\]
The bare pentagon has normalized polynomial $Z_5(t)+2i$ and phase
$\theta_5(t)=\arctan(2/Z_5(t))$.  If $B-A\leq0$, then
$\Theta_H\leq0<\theta_5$.  If $B-A>0$, inequality~(3.9) gives
$R>Z_5(t)(B-A)$ and again $\Theta_H<\theta_5$.  Thus the comparison is strict
for every $t>0$.

The eigenvalues of $C_5$ are $2\cos(2\pi j/5)$, $0\leq j<5$.  Its positive
eigenvalues are $2$ and $(\sqrt5-1)/2$ twice, so
\[
 \splus(C_5)=7-\sqrt5,
 \quad \sminus(C_5)=3+\sqrt5,
 \quad \D(C_5)=4-2\sqrt5=-2\delta.                       \tag{3.10}
\]
Lemma~\ref{lem:coulson} and the strict phase comparison give
$\D(H)>\D(C_5)=-2\delta$.  Finally, $H$ has $a+1$ cyclic blocks, hence
$|E(H)|=|V(H)|+a$, and
\[
 \splus(H)+\sminus(H)=\tr A(H)^2=2|E(H)|.
\]
Averaging this identity with the bound on $\D(H)$ yields
$\sig(H)>a-\delta>0$, proving~(L).
\end{proof}

The proof explicitly covers a pentagon sharing several distinct cuts with the
packing-one triangular part and a pentagon separated from that part by an
arbitrarily branching bridge tree.  No one-interface or common-cut
classification is hidden in Lemma~\ref{lem:L}.

\section{Well-founded actual-bridge induction}\label{sec:bridges}

Fix a $2\Pp$ Voronoi territory $K$.  All its retained triangles have packing
number one, a property inherited by every induced component obtained by
cutting an actual bridge.

A \emph{shared-cut cluster} is a maximal collection of complete cyclic blocks
connected by successive shared cut vertices.  Contract each cluster support
in the cactus block--cut tree.  Between distinct cluster nodes lies a nonempty
chain of actual bridge blocks.  Consider only bridges on the minimal subtree
spanning the cyclic clusters.  Cutting one such bridge gives two connected
induced components, each containing a complete cycle.  Every unused connector
segment, bridge-tree branch, cycle fragment, and off-hull tree remains on the
side containing its actual anchor.  Complete cycles are recomputed in the
physical components after each cut.

\begin{proposition}[Actual-bridge reduction]\label{prop:bridge}
Every $2\Pp$ packing-one territory with more than one shared-cut cluster has
positive surplus, provided the same assertion holds for smaller $2\Pp$
territories.
\end{proposition}

\begin{proof}
Use lexicographic induction on
\[
 \mu(K)=(\text{number of complete cyclic blocks},
          \text{number of shared-cut clusters}).          \tag{4.1}
\]
Only hull bridges whose two sides contain complete cycles are eligible.  If a
cut leaves one terminal side and an active $\T^b\Pp\Pp$ side, the active side
is missing every complete cycle on the terminal side.  Its first coordinate
in~(4.1) is therefore strictly smaller.  Thus recursion cannot repeat the same
instance, and connector-only branches cannot create an infinite or vacuous
step.

For an eligible bridge, the distribution of the two pentagons is exhaustive.
\begin{enumerate}
\item One side has no pentagon.  It is a nonempty triangular packing-one
cactus and has nonnegative surplus, in fact strict by
Lemma~\ref{lem:inputs}\ref{in:tri}.  If the opposite $2\Pp$ side retains a
triangle, it is a smaller induction instance.  If it retains none, it is one
connected $\Pp\Pp$ packet with nonnegative surplus, and the triangular side
supplies strictness.

\item One side is a bare one-pentagon packet.  The exact complete-cycle
profile is
\[
 \Pp+\T^b\Pp,
\]
with $b\geq1$.  Lemma~\ref{lem:inputs}\ref{in:pents} and Lemma~\ref{lem:L}
give
\[
 \sig(K)\geq-\delta+(b-\delta)=b-2\delta>0.              \tag{4.2}
\]
Lemma~\ref{lem:L} applies to the whole second component with its actual
incidence, not to an imagined marked interface.

\item One side is a pure connected $\Pp\Pp$ packet.  Keep it intact, with
surplus at least zero.  The other side is a nonempty triangular packing-one
cactus and is strict.  In particular, the two pentagons are not split into two
adverse singleton charges.

\item Each side contains one pentagon and at least one triangle.  Both are
packing-one one-pentagon packets, so Lemma~\ref{lem:L} makes both positive.
\end{enumerate}
These are exactly the $0+2$ and $1+1$ pentagon distributions and their
symmetric forms.  A side called $\Pp$ or $\Pp\Pp$ has no hidden triangle:
packet selection occurs only after reading complete cycles from the actual
bridge component.  Superadditivity completes every induction step.
\end{proof}

It remains to prove positivity when all complete cycles of $K$ form one
shared-cut cluster.

\section{One-cluster structure and the seven forms}\label{sec:structure}

For one shared-cut cluster, form the bipartite cycle--cut incidence graph $I$:
its nodes are complete cyclic blocks and cut vertices belonging to at least
two such blocks, with containment as incidence.  The graph $I$ is a tree.  An
incidence cycle would produce either a graph cycle traversing several blocks
or two cyclic blocks meeting in more than one vertex, both forbidden in a
cactus.

Let the cluster have $a\geq1$ triangles and two pentagons.  If $a=1$, it is a
$\T\Pp\Pp$ tricyclic cactus and is strict by
Lemma~\ref{lem:inputs}\ref{in:small}; this includes the incidence path
$\Pp-x-\T-y-\Pp$.  Assume henceforth that $a\geq2$.

Choose triangles $T_1,T_2$.  They meet because the triangle-packing number is
one; call their common vertex $x$.  Every third triangle $T$ meets both.  If
$T$ met $T_1$ away from $x$ and $T_2$ away from $x$, the corresponding paths
in $I$ would form an incidence cycle.  If it met one away from $x$ and the
other at $x$, it would share two vertices with that triangle.  Both are
impossible.  Hence every triangle contains $x$.

After deleting the pentagon nodes and cuts used only by them, the triangular
incidence is a star centered at $x$.  A non-$x$ cut belongs to at most one
triangle, since two such triangles would share it as well as $x$.  Call these
non-$x$ vertices \emph{private} triangle ports.

Restore the pentagons.  A pentagon independent of the other can attach to the
triangular star at only one cut; two attachments would close a cycle in $I$.
The cut is either $x$ or a private port of one triangle.  If one pentagon lies
between the other and the star, the two pentagons form a serial pair, whose
first pentagon attaches at $x$ or at a private port.  Therefore, up to
exchanging the pentagons and triangle petals, the possibilities are exactly
\begin{align*}
\mathrm{H1}&:\ \text{both pentagons attach at }x;\\
\mathrm{H2}&:\ \text{both attach at the same private cut }y
                    \text{ of one triangle};\\
\mathrm{H3}&:\ \text{one attaches at }x\text{ and one at a private cut }y;\\
\mathrm{H4}&:\ \text{they attach at private cuts on two distinct triangles};\\
\mathrm{H5}&:\ \text{a serial }\Pp\text{--}\Pp\text{ pair whose first
                    pentagon attaches at }x;\\
\mathrm{H6}&:\ \text{a serial }\Pp\text{--}\Pp\text{ pair whose first
                    pentagon attaches at a private cut }y;\\
\mathrm{H7}&:\ \text{they attach at the two distinct private cuts }y,z
                    \text{ of one triangle}.             \tag{5.1}
\end{align*}
This list is forced by the tree $I$: one private port gives H2, two ports on
one petal give H7, ports on distinct petals give H4, mixed central/private
ports give H3, and dependence of the pentagon nodes gives exactly the two
serial cases H5--H6.  There is no triangle hidden beyond a pentagonal port,
because every triangle is already a petal at $x$.

\section{Physical ownership for H1--H7}\label{sec:H}

We first record the physical rule used in the non-hinge rows.  When vertices
of a router cycle are partitioned into consecutive intervals, the entire
incidence component at an occupied port follows the interval containing that
port.  Every connector remnant and every arbitrary tree follows its actual
anchor.  Edges joining different owners disappear only when induced
subgraphs are formed.  Thus owners are connected, induced, disjoint, and
exhaustive; a proper cycle interval is a path and creates no complete cycle.

\begin{proposition}[The seven terminal forms]\label{prop:H}
Every one-cluster packing-one $\T^a\Pp\Pp$ cactus has positive surplus.
More precisely, for $a\geq2$ the exact terminal packets and margins are:
\begin{center}
\begin{tabular}{@{}cll@{}}
\toprule
Type & Physical terminal packets & Surplus bound\\
\midrule
H1 & common-cut $\T^a\Pp\Pp$
   & $>a+1-4/(3\sqrt{13})$\\
H2 & rooted common-cut $\Pp\Pp$ hinge entered at $y$
   & $>0$\\
H3 & $\Pp+\T^{a-1}\Pp$ via a $1+2$ triangle split
   & $>(a-1)-2\delta$\\
H4 & $\Pp+\T^{a-1}\Pp$ via a pentagon-bearing petal
   & $>(a-1)-2\delta$\\
H5 & rooted $\Pp\Pp$ hinge entered at $x$
   & $>0$\\
H6 & rooted $\Pp\Pp$ hinge entered at $y$
   & $>0$\\
H7 & $\Pp+\T^{a-1}\Pp$ via a valid $1+2$ port split
   & $>(a-1)-2\delta$\\
\bottomrule
\end{tabular}
\end{center}
\end{proposition}

\begin{proof}
H1 is exactly Lemma~\ref{lem:inputs}\ref{in:common}.  In H2 the two
pentagons share their cut and the nonempty triangular star enters that rooted
hinge at $y$ through one interface.  H5 and H6 are the same hinge dependency,
with entry at $x$ and $y$, respectively.  These three applications include all
finite trees at the hinge, triangular side, and interface; no bare-core
replacement is made.

For H3, split the marked router triangle into the singleton interval $\{y\}$
and the complementary two-vertex interval.  The first owner receives $y$, all
four private vertices of the pentagon incident at $y$, and every remnant or
tree anchored in that incidence branch.  It therefore retains the whole
pentagon: ``singleton $y$'' refers only to the interval of the router triangle,
not to a singleton-vertex final territory.  The complementary owner receives
the other two triangle vertices, every remaining incidence branch, and every
remaining vertex.  Its complementary triangle edge keeps it connected through
$x$.  The router triangle is destroyed; this owner retains exactly $a-1$
triangles and the other complete pentagon.  Hence the physical packet profile
is $\Pp+\T^{a-1}\Pp$.

The H4 split is identical on either pentagon-bearing petal: choose its occupied
private port as the singleton interval and give the other two triangle
vertices, including $x$, to the complementary owner.  The other pentagon's
petal and all remaining petals stay connected through $x$, again giving
$\Pp+\T^{a-1}\Pp$.

For H7, the router triangle has three distinct ports $x,y,z$.  Choose either
occupied private port, say $y$, as the singleton interval.  Its entire
pentagon branch receives $y$.  The other two consecutive triangle vertices
form one interval containing $x$ and $z$; they receive the other pentagon
branch and all triangular petals at $x$.  This is a valid physical $1+2$
split, destroys the router triangle, and has the same
$\Pp+\T^{a-1}\Pp$ profile.  Each shared cut is assigned once, and every
incidence branch follows its actual port.  This is an ownership construction,
not an abstract demand allocation.

In H3, H4, and H7, the one-pentagon owner retains a subfamily of the original
packing-one triangles, so Lemma~\ref{lem:L} applies.  Charging the bare
pentagon by Lemma~\ref{lem:inputs}\ref{in:pents} gives
\[
 \sig(K)>(a-1)-2\delta
 \geq1-2\delta=5-2\sqrt5>0.                              \tag{6.1}
\]
This proves all seven rows.
\end{proof}

The H7 row is asserted for the unmarked one-cluster object classified in
Section~\ref{sec:structure}: its private occupied port is available, the other
two vertices are one consecutive interval containing $x$ and the other port,
and every branch follows its actual port.  It is not a rule that may be copied
blindly into an externally marked recursive state.  If an external entry would
disconnect an owner or assign a port twice, that opening is unavailable.  At
total ranks four and five the established $\T^2\Pp\Pp$ and
$\T^3\Pp\Pp$ terminals are used.  At larger marked rank, a leaf-pentagon
opening is valid only after checking physically that deleting its occupied cut
leaves its four private vertices, remnants, and attachments in one connected
acyclic owner and leaves a connected packing-one $\T^a\Pp$ complement.  When
that predicate holds, its ledger is
\[
 -1+(a-\delta)>0
\]
by Lemma~\ref{lem:L}; when it fails, no formal H7 opening is claimed.  The main
proof needs only the unmarked form already proved in Proposition~\ref{prop:H}.

\section{Completion, dependencies, and scope}\label{sec:dependencies}

\begin{proof}[Proof of Theorem~\ref{thm:main}]
First apply maximum-triangle Voronoi to $G$.  Make no bridge cut beforehand.
Lemma~\ref{lem:voronoi} says that every territory contains its selected
triangle and has triangle-packing number one.  The $0\Pp$, $1\Pp$, and
$2\Pp$ classification is exhaustive.  The first type is strict by
Lemma~\ref{lem:inputs}\ref{in:tri}, and the second is positive by the
self-contained Lemma~\ref{lem:L}.  For the third type, induct on the
well-founded measure~(4.1).  Proposition~\ref{prop:bridge} handles every
multiple-cluster step, while Proposition~\ref{prop:H} is the one-cluster base.
Superadditivity over the nonempty induced territories gives
\[
 \sig(G)\geq\sum_i\sig(G_i)>0.
\]
Equivalently, $\splus(G)>|V(G)|$.
\end{proof}

The proof depends on the following previously proved results, in addition to
the elementary superadditivity and the complete proof of Lemma~\ref{lem:L}
given here:
\begin{itemize}
\item the packing-one triangular phase estimate and triangular Voronoi
territories from the all-rank one-hostile-cycle manuscript;
\item the exact common-cut $\T^b\Pp\Pp$ Schur--Sachs inequality for H1;
\item the rooted two-pentagon hinge theorem, including its separate
common-root deletion identities, for H2, H5, and H6;
\item the complete tricyclic theorem for $a=1$, and the proved rank-four and
rank-five terminals for externally marked low-rank fallback;
\item the physical consecutive-interval router lemma for exact ownership of
ports, cuts, remnants, and attached trees.
\end{itemize}
These dependencies are the repository proof objects
\cite{Hostile,CommonCut,Hinge,Tricyclic,Tetracyclic,Pentacyclic,Router}.
In particular, no unrestricted global estimate
$\sig(\T^a\Pp)>a-\delta$ is assumed.  Every quantitative one-pentagon
invocation occurs inside a Voronoi territory or one of its induced
bridge/router descendants, where the retained triangles still have packing
number one and Lemma~\ref{lem:L} applies.

\paragraph{Scope and nonclaim.}
The theorem covers finite simple connected cacti with exactly $r\geq1$
triangles and exactly two pentagons, arbitrary shared cuts, arbitrary bridge
lengths and branching, and arbitrary finite forest attachments.  Split
pentagon fragments remain forests with their physical owners.  The proof does
not cover three pentagons, non-cactus block intersections, or a merely maximal
triangle packing.  It establishes the stated family only.  It is not a proof
of the universal Akbari--Kumar--Mohar--Pragada--Zhang positive-square-energy
conjecture for all graphs, nor even a claim for every cactus outside this block
profile.

\section*{AI disclosure}

This manuscript was generated and assembled by an AI coding and mathematical
reasoning system from the accepted proof note and the cited repository proof
objects.  No human authorship, human review, or human verification is claimed.
The result should be assessed from the displayed argument and its explicit
dependencies rather than from asserted authority.

\begin{thebibliography}{99}
\bibitem{AKMPZ}
S.~Akbari, H.~Kumar, B.~Mohar, S.~Pragada, and X.~Zhang,
\emph{Refinement of a conjecture on positive square energy of graphs},
arXiv:2506.07264v1, 2025.

\bibitem{Hostile}
Companion manuscript, \emph{Cacti with Arbitrarily Many Triangles and One
Hostile Cycle}, 2026;
\path{all-rank-triangle-hostile-cacti/paper.tex}.

\bibitem{CommonCut}
\emph{Common-cut residual bouquets: an exact rooted Schur--Sachs inequality},
26 July 2026;
\path{research/common-cut-bouquet-rooted-schur-2026-07-26.md}.

\bibitem{Hinge}
\emph{Rooted two-pentagon hinge theorem}, 28 July 2026;
\path{research/rooted-two-pentagon-hinge-theorem-2026-07-28.md}.

\bibitem{Tricyclic}
Companion manuscript, \emph{Positive Square Energy of Every Tricyclic Cactus},
2026; \path{all-tricyclic-cacti/paper.tex}.

\bibitem{Tetracyclic}
Companion manuscript, \emph{Positive Square Energy of Every Tetracyclic
Cactus}, 2026; \path{all-tetracyclic-cacti/paper.tex}.

\bibitem{Pentacyclic}
Companion manuscript, \emph{All Pentacyclic Cacti}, 2026;
\path{all-pentacyclic-cacti/paper.tex}.

\bibitem{Router}
\emph{Physical $C_5$ interval-router lemma}, 29 July 2026;
\path{research/physical-c5-interval-router-lemma-2026-07-29.md}.
\end{thebibliography}

\end{document}
